% Generated draft: requires mathematical and source review.
\documentclass[11pt]{article}
\usepackage[margin=25mm]{geometry}
\usepackage{fontspec}
\setmainfont{Latin Modern Roman}
\usepackage{amsmath,amssymb,mathtools,unicode-math}
\setmathfont{Latin Modern Math}
\usepackage{enumitem,needspace,xurl,hyperref}
\setlength{\parindent}{0pt}
\setlength{\parskip}{5pt}
\setlength{\emergencystretch}{3em}
\newcommand{\N}{\mathbb{N}}
\newcommand{\Z}{\mathbb{Z}}
\newcommand{\Q}{\mathbb{Q}}
\newcommand{\R}{\mathbb{R}}
\newcommand{\C}{\mathbb{C}}
\newcommand{\F}{\mathbb{F}}
\newcommand{\A}{\mathbb{A}}
\newcommand{\PP}{\mathbb{P}}
\newcommand{\E}{\mathbb{E}}
\newcommand{\eps}{\varepsilon}
\newcommand{\dd}{\,\mathrm{d}}
\begin{document}
\section*{UnsolvedMath ID: 11000263}
\textbf{Source problem code:} catalogue entry
\subsection{Definitions and mathematical statement}
The algebraic setup is supplied in the preceding local record 11000262. Let $B_n$ be the braid group and $RB_n$ its group algebra over the source's coefficient ring $R$, with parameter $q$. Write $\sigma_{i_1\cdots i_k}=\sigma_{i_1}\cdots\sigma_{i_k}$, and retain the source's barred-word notation $\bar\sigma_{i_1\cdots i_k}$ for its inverse-word convention. The extraction does not clearly distinguish inversion of the whole word from replacing letters by inverses in the written order; this matters for words of length greater than one.

The source defines
\[
 X_2=q\sigma_1^{-1}+(1-q)1-\sigma_1,\qquad
 X_k=(q^{k-1}\bar\sigma_{(k-1)\cdots1}-\sigma_{1\cdots(k-1)})X_{k-1}\quad(k\ge3).
\]
It defines $Z_n$ as a quotient of $RB_n$ by the relations beginning with
\[
 (q\sigma_2^{-1}+1-q-\sigma_2)X_2
 =(q\bar\sigma_{21}-\sigma_{12})X_2,
\]
and continuing in the pattern
\[
 (q^{j-1}\bar\sigma_{j\cdots2}-\sigma_{2\cdots j})X_j
 =(q^{j-1}\bar\sigma_{j\cdots1}-\sigma_{1\cdots j})X_j.
\]
The written list ends with $j=n$. Under the standard $B_n$ convention only $\sigma_1,\ldots,\sigma_{n-1}$ exist, so this endpoint requires correction or a different braid indexing convention. No endpoint correction is silently chosen here. The ring $R$ and its parameters also are not specified in the extracted setup.

\textbf{Question 6.} Does the element $X_3$ vanish in $Z_n$? Once the coefficient, inverse-word, and index conventions are fixed, this is the ideal-membership question
\[
 X_3\in I_n\subset RB_n,
\]
where $I_n$ is the two-sided ideal generated by the defining relations; $n$ must be in a range where $X_3$ and the relations are defined. It is not asking whether an independent generator named $X_3$ may be declared zero.

The source suggests that additional relations may be needed, for example
\[
 \sigma_1X_2=tX_2,
\]
or a more general relation. The parameter $t$ and its relation to $R$ require specification. This suggestion is not already imposed in the question. The preceding record says the Iwahori--Hecke algebra $H_n(1,-q)$ is obtained by imposing $X_2=0$, and the Birman--Wenzl--Murakami quotient by imposing both $X_3=0$ and $\sigma_1X_2=tX_2$. Thus the question probes whether the proposed $Z_n$ genuinely retains more information than that quotient; a quotient relation cannot itself prove $X_3=0$ in the original algebra.
\subsection{Short English statement}
Determine whether $X_3$ is already zero in Bigelow's proposed braid-group algebra quotient, before adding the relations that recover the Birman--Wenzl--Murakami algebra. The extracted definition has inverse-word, coefficient-ring, and generator-index ambiguities.
\subsection{Sources}
\begin{itemize}
\item Local UnsolvedMath registry, permanent ID 11000263, source code catalogue entry; no separate proposer or proposal year supplied.
\item Farb, edited volume \emph{Problems on Mapping Class Groups and Related Topics}, Question 6, PDF page 323, in source text bearing S. Bigelow's chapter attribution; \url{https://www.math.uchicago.edu/~farb/papers/mcgbook.pdf}. The registry records PDF-font-structure extraction and access on 2026-07-29; its inventory distinguishes 329 labeled headings by source page.
\item Adjacent local ID 11000262 contains the $X_k$ recursion, $Z_n$ relations, and the Hecke/Birman--Wenzl--Murakami quotient descriptions. Its homological motivation is expressly described by the source as unclear and unpublished at that time.
\end{itemize}
\end{document}
